\documentclass[11pt,a4paper]{article}

% PACKAGES
\usepackage[utf8]{inputenc}
\usepackage[T1]{fontenc}
\usepackage[slovak]{babel}
\usepackage{amsmath, amssymb}
\usepackage{amsthm} 
\usepackage{graphicx}
\usepackage{enumerate}
\usepackage{hyperref}

\usepackage{tikz}
\usetikzlibrary{tikzmark}
\usetikzlibrary{arrows.meta}
\usetikzlibrary{calc}
\usetikzlibrary{shapes.geometric} 
\usetikzlibrary{positioning}      
\usetikzlibrary{cd}
\usetikzlibrary{babel}
\usetikzlibrary{decorations.pathreplacing}
\usetikzlibrary{angles} 
\usetikzlibrary{quotes}

\usepackage[all,pdf,2cell]{xy}
\usepackage{soul}
\usepackage{framed}
\usepackage{array}
\usepackage{todonotes}
\usepackage{nicematrix}

\NiceMatrixOptions
  { pgf-node-code = \pgfsetfillcolor{white} \pgfusepathqfill }

\pgfset{nicematrix/cell-node/.append style = { inner sep = 1pt } }

% THEOREM ENVIRONMENTS SETUP
\theoremstyle{definition} 
\newtheorem{definicia}{Definícia}[section]
\newtheorem{priklad}[definicia]{Príklad}

\theoremstyle{plain} 
\newtheorem{tvrdenie}[definicia]{Tvrdenie}
\newtheorem{veta}[definicia]{Veta}
\newtheorem{dosledok}[definicia]{Dôsledok}

\theoremstyle{remark} 
\newtheorem*{question}{Otázka}

% Macros
\newcommand{\id}{\mathrm{id}}
\newcommand{\Nat}{\mathbb{N}}
\newcommand{\Set}{\mathbf{Set}}
\newcommand{\R}{\mathbb{R}}
\newcommand{\C}{\mathbb{C}}
\newcommand{\ora}[1]{\overrightarrow{#1}}
\newcommand{\Lo}{\mathrm{Lo}} 

\let\Im\relax
\let\Re\relax
\DeclareMathOperator{\Im}{Im}
\DeclareMathOperator{\Re}{Re}
\DeclareMathOperator{\Ker}{Ker}
\DeclareMathOperator{\proj}{proj}

\newcommand\restr[2]{{%
  \left.\kern-\nulldelimiterspace 
  #1 
  \littletaller 
  \right|_{#2} 
  }}
\newcommand{\littletaller}{\mathchoice{\vphantom{\big|}}{}{}{}}

\begin{document}

\section*{DIAGONALIZÁCIA, MOTIVÁCIA}

Najskôr istá zmena: lineárne zobrazenia budem značiť veľkým písmenom lebo „tak sa to robí“. V skutočnosti tam nie je prevládajúca konvencia, mení sa to situačne.

\setcounter{section}{3} % Iba kvôli číslovaniu 3.01
\setcounter{definicia}{0}

\begin{definicia}
\ul{Lineárna transformácia} je lineárne zobrazenie vektorového priestoru do seba:
\end{definicia}

\vspace{0.3cm}
\begin{center}
\begin{tikzpicture}
\node (eq) at (0,0) {$A \colon V \longrightarrow V \quad \square$};
\node[red, align=center] (desc) at (0,-1.3) {lineárna transformácia\\priestoru $V$ / na priestore $V$.};
\draw[->, red, thick] (desc) -- (0,-0.3);
\end{tikzpicture}
\end{center}
\vspace{0.3cm}

Ak $V$ je konečnorozmerný, každá matica $A$ je štvorcová, samozrejme.

Cieľom tejto, aj ďalších prednášok je hľadať bázu $X$ tak, aby
\[ [A]_{XX} \text{ bola čo najjednoduchšia} \]
podľa možnosti diagonálna.

\begin{definicia}
Lineárna transformácia $A \colon V \to V$ na konečnorozmernom vektorovom priestore $V$ je \ul{diagonalizovateľná}, ak existuje báza $X$ vektorového priestoru $V$ taká, že $[A]_{XX}$ je diagonálna matica.
\end{definicia}

Prirodzené otázky:
\begin{itemize}
    \item Je každá lin. transformácia diagonalizovateľná? [nie]
    \item Ktoré sú? [uvidíme]
    \item Ak nejaká transformácia má dve diagonálne matice $[A]_{XX}$, $[A]_{YY}$ v dvoch rôznych bázach, musia byť rovnaké? [skoro]
\end{itemize}

Jedna z vecí, ktoré idú dobre robiť s diagonálnymi maticami: ľahko sa násobia medzi sebou, medziiným sa ľahko umocňujú: (iba príklad)

\vspace{0.5cm}
\noindent \textbf{Motivácia: Fibonacciho postupnosť}

\begin{minipage}{0.65\textwidth}
\[ F_0 = 1 \quad F_1 = 1 \quad F_{n+2} = F_{n} + F_{n+1} \quad n \in \Nat \]
\end{minipage}
\begin{minipage}{0.3\textwidth}
\begin{center}
$1 \ 1 \ 2 \ 3 \ 5 \ 8$\\
\small $0 \ 1 \ 2 \ 3$
\end{center}
\end{minipage}

\[
\begin{pmatrix} 1 & 1 \\ 1 & 0 \end{pmatrix} 
\begin{pmatrix} F_{n+1} \\ F_n \end{pmatrix} 
= \begin{pmatrix} F_n + F_{n+1} \\ F_{n+1} \end{pmatrix}
\]

\vspace{0.3cm}
\noindent $A \colon \R^2 \to \R^2$

\[
A \begin{pmatrix} F_1 \\ F_0 \end{pmatrix} = \begin{pmatrix} F_2 \\ F_1 \end{pmatrix} \qquad 
A \begin{pmatrix} F_2 \\ F_1 \end{pmatrix} = \begin{pmatrix} F_3 \\ F_2 \end{pmatrix}
\]

\[
A^2 \begin{pmatrix} F_1 \\ F_0 \end{pmatrix} = \begin{pmatrix} F_3 \\ F_2 \end{pmatrix} \quad \dots \quad 
\underbrace{A^n}_{\color{red}\begin{matrix} \downarrow \\ \text{toto je škaredá} \\[-0.1cm] \text{matica} \end{matrix}} 
\begin{pmatrix} F_1 \\ F_0 \end{pmatrix} = \begin{pmatrix} F_{n+1} \\ F_n \end{pmatrix}
\]

\vspace{0.5cm}
\begin{center}
\begin{tikzpicture}
\node (eq) at (0,0) {$\displaystyle [A]^n_{XX} \left[\begin{pmatrix} 1 \\ 0 \end{pmatrix}\right]_X = \left[ \begin{pmatrix} F_{n+1} \\ F_n \end{pmatrix} \right]_X$};
\node[align=left, text width=10.5cm] (desc) at (1,-1.5) {ak nájdeme takú bázu $X$ vektorového priestoru $\R^2$, že $[A]_{XX}$ je diagonálna matica, vieme toto ľahko zrátať.};
\draw[->, thick] ($(eq.south west) + (1.2,0.2)$) to[out=-90, in=180] ($(desc.west) + (0,0)$);
\end{tikzpicture}
\end{center}
\vspace{0.5cm}

\section*{ZMENA BÁZY, PODOBNOSŤ MATÍC}

Povedzme, že máme danú lineárnu transformáciu $A \colon V \to V$ na konečnorozmernom vektorovom priestore $V$, pomocou matice $[A]_{YY}$. Ak máme inú bázu $X$, ako nájdeme $[A]_{XX}$?

\begin{align*}
[A]_{XX} &= [\id_V \circ A \circ \id_V]_{XX} = \\[0.3cm]
&= [\id_V]_{XY} \circ [A]_{YY} \circ \underbrace{[\id_V]_{YX}}_{\color{red}\begin{matrix} \downarrow \\ \text{ako vyzerá táto matica?} \end{matrix}}
\end{align*}

Ak $X = (\ora{x_1}, \dots, \ora{x_n})$

\begin{align*}
[\id_V]_{YX} &= \begin{pmatrix} 
[\id_V(\ora{x_1})]_Y & \dots & [\id_V(\ora{x_n})]_Y 
\end{pmatrix} = \\[0.3cm]
&= \begin{pmatrix} 
[\ora{x_1}]_Y & \dots & [\ora{x_n}]_Y 
\end{pmatrix}
\end{align*}

ak $V = K^n$ a $Y$ je kanonická báza
\begin{center}
\begin{tikzpicture}
\node (eq1) at (0,0) {$Y = E = (\ora{e_1}, \dots, \ora{e_n})$};
\node (text1) at (0,-1) {jednotkové vektory $K^n$};
\draw[->] (eq1) -- (text1);

\node (eq2) at (0,-2.5) {$[\ora{x_i}]_Y = \ora{x_i}$};
\node (text2) at (0,-3.5) {ako stĺpec};
\draw[->] (eq2) -- (text2);
\end{tikzpicture}
\end{center}

\vspace{0.3cm}
\noindent Aký je vzťah $[\id_V]_{XY}, [\id_V]_{YX}$?
\begin{align*}
[\id_V]_{XY} [\id_V]_{YX} &= [\id_V \circ \id_V]_{XX} \\
&= I_n \quad \textbf{!}
\end{align*}

Teda sú navzájom inverzné.

Pr.: Lin. transf: osová súmernosť dla osi I. kvadrantu $A$:
\[
[A]_{EE} = \begin{pmatrix} 0 & 1 \\ 1 & 0 \end{pmatrix} \qquad \text{AKA „prehodenie zložiek“}
\]

Vezmime $X = ((1,1), (-1,1))$
\[
\begin{pmatrix} 1 & -1 \\ 1 & 1 \end{pmatrix}^{-1} = \begin{pmatrix} 1/2 & 1/2 \\ -1/2 & 1/2 \end{pmatrix}
\]

\begin{align*}
[A]_{XX} &= \begin{pmatrix} 1/2 & 1/2 \\ -1/2 & 1/2 \end{pmatrix} 
\begin{pmatrix} 0 & 1 \\ 1 & 0 \end{pmatrix} 
\begin{pmatrix} 1 & -1 \\ 1 & 1 \end{pmatrix} = \\[0.2cm]
&= \begin{pmatrix} 1/2 & 1/2 \\ 1/2 & -1/2 \end{pmatrix} 
\begin{pmatrix} 1 & -1 \\ 1 & 1 \end{pmatrix} = 
\begin{array}[t]{c}
\begin{pmatrix} 1 & 0 \\ 0 & -1 \end{pmatrix} \\
\downarrow \\
\text{diagonálna!}
\end{array}
\end{align*}

\vspace{0.3cm}
Toto motivuje nasledujúcu definíciu:

\begin{definicia}
Dve štvorcové matice rovnakého rozmeru $A, B$ sú \ul{podobné}, ak existuje regulárna matica $P$ taká, že
\[ B = P^{-1} A P \]
Značíme $A \simeq B$
\end{definicia}

\begin{tvrdenie}
Ak $A, B, C$ sú štvorcové matice $n \times n$
\begin{enumerate}[a)]
    \item $A \simeq A$
    \item Ak $A \simeq B$, potom $B \simeq A$
    \item Ak $A \simeq B$ a $B \simeq C$, potom $A \simeq C$
\end{enumerate}
\end{tvrdenie}

\begin{proof}
\hfill
\begin{enumerate}[a)]
    \item $A = I^{-1} A I$
    \item Ak $B = P^{-1} A P$, potom
    \begin{align*}
    P B &= P P^{-1} A P = I A P = A P \\
    P B P^{-1} &= A P P^{-1} = A I = A
    \end{align*}
    \[ (P^{-1})^{-1} B P^{-1} = A \implies B \simeq A \]
    \item $B = P^{-1} A P \qquad C = Q^{-1} B Q$
    \begin{align*}
    C &= Q^{-1} P^{-1} A P Q = \\
    &= (P Q)^{-1} A (P Q) \\[0.2cm]
    &\Downarrow \\
    A &\simeq C \qedhere
    \end{align*}
\end{enumerate}
\end{proof}

\vspace{0.5cm}
\noindent Pohľad na vec ako „diagonalizáciu matice“

\begin{center}
\begin{tikzpicture}
% Circle representing the set of matrices
\draw[thick] (0,0) circle (2cm);

% Slices representing equivalence classes
\begin{scope}
\clip (0,0) circle (2cm);
\draw[thick] (-3, -0.8) -- (3, 0.8);
\draw[thick] (-3, 0.5) -- (3, 2);
\draw[thick] (-3, -2) -- (3, -0.5);
\draw[thick] (-0.5, -3) -- (1.5, 3);
\draw[thick] (-2, -3) -- (0, 3);
\draw[thick] (1, -3) -- (3, 3);
\end{scope}

% \simeq symbol in one class
\node at (-0.6, 0.1) {\large $\simeq$};

% Representative matrices as dots in their classes
\filldraw (0.4, 0.5) circle (1.5pt);
\filldraw (0.8, -0.1) circle (1.5pt);
\filldraw (0.5, -0.6) circle (1.5pt);

% Arrows to "diagonálne"
\node[right] (diag) at (2.5, 0) {diagonálne};
\draw[->, thick, shorten <=3pt] (0.4, 0.5) to[out=0, in=170] (diag.west);
\draw[->, thick, shorten <=3pt] (0.8, -0.1) -- (diag.west);
\draw[->, thick, shorten <=3pt] (0.5, -0.6) to[out=0, in=190] (diag.west);

% Text K^{nxn}
\node[above right] at (1.5, 1.5) {$K^{n \times n}$ matice $n \times n$};

% Bottom arrow and text
\draw[->, thick] (-0.5, -2.1) -- (-0.5, -2.8);
\node[below] at (-0.5, -2.8) {rozklad $K^{n \times n}$ podľa podobnosti};
\end{tikzpicture}
\end{center}

\vspace{0.5cm}
Hľadáme podobnú, diagonálnu maticu.

\end{document}
